% The six blocks, in the order index.html places them.
%
% Each macro carries the SUMMARY view of its page verbatim: the detail views
% are the "Details" pages the print poster does without, so nothing from them
% appears here. Columns 2 and 3 hold more text than their slots would take at
% full size, so those two blocks set a fit scale; the rest run at 1.

% =========================================================== 1. The puzzle ==
\newcommand{\blockOne}{%
    \setfit{0.95}%
    \kicker{The puzzle}
    \htitle{Automation exposure is widespread, yet experimental evidence on
        policy demand remains inconclusive.}

    \statrow{%
        \stat{bad}{\statw{2}}
            {High-risk occupations}{47\%}{of U.S. employment (Frey \& Osborne 2017)}%
        \statgap
        \stat{warn}{\statw{2}}
            {Public concern}{62\%}{major concern about AI at work (Pew 2023)}}
    \vgap

    \hsub{Prior evidence}
    \begin{plainlist}
        \plainitem{Zhang \& Hou (2022); Magistro et al. (2024).}{Personalized
            risk information produced a null average effect.}
        \plainitem{Gallego et al. (2022).}{A different treatment and sample
            produced positive effects.}
        \plainitem{Arntz et al. (2024).}{Opposing subgroup effects cancelled in
            the aggregate.}
    \end{plainlist}

    \callout{accent}{The diagnosis}{Heterogeneity washes out in the aggregate.
        Weak average effects may reflect composition rather than stable
        preferences.}%
}

% ================================================ 2. Theory and hypotheses ==
\newcommand{\blockTwo}{%
    \setfit{1.2}%
    \kicker{Theory \& hypotheses}
    \htitle{When do citizens demand state protection in response to expert
        warnings about automation risk?}
    \lede{The Conditional Activation Model answers with two sequential gates.}

    \begin{plainlist}
        \plainitem{H1 · Trust gate.}{Expert warnings raise policy demand only
            among respondents who credit the source. Among respondents who do
            not, the effect is zero.}
        \plainitem{H2 · Joint activation.}{Among respondents who credit the
            source, the effect is largest where trait anxiety and
            own-occupation risk are both high. Neither anxiety nor risk alone
            produces it.}
        \plainitem{H3 · Composition.}{The average effect understates the
            effect in the responding subgroup, because it averages over
            respondents for whom the gates did not clear.}
    \end{plainlist}%
}

% ======================================================= 3. Research design ==
\newcommand{\blockThree}{%
    \setfit{0.865}%
    \kicker{Research design}
    \htitle{Pre-registered between-subjects experiment, N = 901 U.S. workers.}
    \lede{A three-step belief-correction protocol whose intensity is
        individualized: every respondent is shown the expert estimate for their
        \emph{own} occupation. 475 treated, 426 control.}

    \begin{stepslist}
        \stepitem{1}{Explain how experts judge automation risk.}%
            {A neutral briefing: tasks that are routine, clearly defined and
             repeated, with plenty of recorded examples, are easier to
             automate. Illustrated with two contrasting jobs, and with what
             raises risk (fully online work, abundant training data) against
             what lowers it (in-person dexterity, reliance on trust or
             empathy). Comprehension is then tested: 93.5\% pass.}
        \stepitem{2}{Elicit the respondent's own belief.}%
            {Every U.S. worker is described as standing in one line ordered by
             automation risk, from position 0 (safest) to position 100
             (riskiest), with 50 as the midpoint. Respondents place their own
             job on that line with a slider, then report how far off they think
             they might be.}
        \stepitem{3}{Reveal the expert estimate and display the gap.}%
            {The source is named: Frey and Osborne at the University of Oxford,
             with their method described. The respondent's own placement and
             the expert position for their occupation are then shown side by
             side, with the difference stated in plain terms.}
    \end{stepslist}

    \callout{accent}{Why the intensity varies}{The treatment is the same
        procedure for everyone, but its content is personalized: the size and
        direction of the gap between belief and expert estimate is whatever
        that respondent's own occupation and prior produce.}

    \hsub{Measures}
    {\tablesetup
    \begin{tabularx}{\linewidth}{@{}Y{0.80}Y{1.62}Y{0.58}@{}}
        \thcell{Construct} & \thcell{Operationalization} & \thcell{Role} \\
        \headrule
        Policy demand (η)    & 5 items, 7-pt (ω = 0.83), graded response model     & Outcome \\ \hline
        Expert trust (S)     & Single 5-pt item, post-treatment, treated arm only  & Gate 1  \\ \hline
        Trait anxiety (A)    & 4-item scale, pre-treatment (ω = 0.92)              & Gate 2  \\ \hline
        Automation risk (R)  & Own-occupation expert percentile, pre-treatment     & Gate 2  \\
    \end{tabularx}\par}
    \vgap

    \callout{accent}{Design property}{Both moderators are measured before
        treatment, so the activation corner is a fixed partition of
        pre-treatment characteristics rather than an outcome of the
        manipulation.}

    \hsub{Identification}
    \para{Expert trust is observed only in the treated arm. A pre-treatment
        proxy carries it into both: a random forest on 15 pre-treatment terms
        predicts the trust rating, a logit on that proxy gives the principal
        score p\textsubscript{i}, and arm means are Hájek-weighted by
        p\textsubscript{i} for the higher-trust stratum and by
        1 − p\textsubscript{i} for the lower-trust one. Identified under
        randomization plus principal ignorability.}%
}

% ============================================================== 4. Results ==
\newcommand{\blockFour}{%
    \setfit{1}%
    \kicker{Results}
    \htitle{Trust gates absorption: the effect among people who trust the
        source is roughly three times the effect among those who trust it
        less.}

    \posterfig{../img/ate-items.png}{The average effect and its per-item
        marginal effects. No single policy drives the composite; support rises
        across the battery, most for income support and labor protections.}

    \posterfig{../img/results-cate.png}{By trust in the expert source. Thick
        and thin bars are the 84\% and 95\% bootstrap intervals; the dotted
        line is the effect for everyone, unweighted.}

    \hsub{Among people who trust the source, the effect concentrates in the
        high-anxiety × high-risk corner}
    {\tablesetup
    \begin{tabularx}{\linewidth}{@{}Y{1.02}Z{0.66}Z{1.20}Z{1.12}@{}}
        \thcell{Trust in source} & \thcell{Outside} & \thcell{Extra in corner}
            & \thcell{In corner} \\
        \headrule
        \hlrow \fBold Higher & \fBold 0.288
            & \fBold 0.397\newline{}[0.083, 0.690]
            & \fBold 0.685\newline{}[0.364, 0.980] \\ \hline
        Lower & 0.092
            & 0.194\newline{}[−0.056, 0.466]
            & 0.285\newline{}[0.041, 0.552] \\
    \end{tabularx}\par}
    \vgap

    \callout{accent}{No corner without the first gate}{Among those who trust
        the source, the corner effect is 2.4× the effect outside it. Among
        those who trust it less, the corner adds nothing detectable. The
        conjunction only pays off for people who took the signal on board.}%
}

% ========================================= 5. Power, robustness, sensitivity ==
\newcommand{\blockFive}{%
    \setfit{1}%
    \kicker{Power, robustness \& sensitivity}
    \htitle{Two of the three headline quantities are well powered; the third
        is not.}

    \statrow{%
        \stat{good}{\statw{3}}{Effect, higher trust}{98\%}{MDE 0.28}%
        \statgap
        \stat{good}{\statw{3}}{Corner, higher trust}{98\%}{MDE 0.44}%
        \statgap
        \stat{warn}{\statw{3}}{Difference between them}{62\%}{below its MDE of 0.33}}
    \vgap

    \para{The higher-trust results are the ones this design can carry.}%
}

% ============================================ 6. Takeaways and implications ==
\newcommand{\blockSix}{%
    \setfit{1.2}%
    \kicker{Takeaways \& implications}
    \htitle{Aggregate effects understate the subgroup that responds.}

    \callout{accent}{Why prior results conflict}{Earlier experiments averaged
        over people who never absorbed the signal. Averaging that way produces a
        weak effect even when part of the sample moves a lot. The nulls in this
        literature fit that pattern.}

    \callout{accent}{Trust decides who responds}{Respondents who trust the
        expert source move about three times as much as those who do not.
        Personal exposure matters only for them. Among trusters the anxiety ×
        risk corner runs 2.4× the effect outside it. Among everyone else it does
        nothing measurable.}

    \callout{warn}{What the next study needs}{This design measured trust after
        treatment, in the treated arm only, so the estimates rest on principal
        ignorability. Measuring trust first removes that assumption. The design
        also fits climate, vaccines and trade.}%
}
